\chapter{What language models contribute to causal knowledge}
\label{la:frontier}

\section{Four different questions about cause and effect}

A banker asking why deposit balances fell can be asking for an explanation of
a recorded event, an estimate of a feasible rate intervention, a prediction of
next month's balance, or a diagnosis of an observation error. A fluent answer
can move between these objects without announcing the change. Evaluating assistance therefore starts by distinguishing the task: a proposed scientific question, evidence about a causal model,
a formal calculation conditional on that model, or a recommendation under a
specified loss. Success on one task does not establish success on the others.

The early frontier argument proposes using language models to organize domain
knowledge, draft graphs, critique assumptions and suggest robustness analyses.
It also records the ability to reproduce recognizable public data and
distinguishes transforming known information from discovering a new causal
structure.\cite{v12_la_2305_00050} This distinction is particularly consequential
in banking. A familiar statement about interest rates can be useful prior
knowledge, but the effect of a particular bank's deposit offer depends on its
population, competitors, eligibility rules and information clock. A remembered
relationship does not determine that intervention effect.

Formal benchmarks introduce a second object. CLadder constructs questions from
known causal structures with different verbalizations. Its results concern
answers to those generated questions, including historically specified models
and prompting methods.\cite{v12_la_2312_04350} Such a benchmark can reveal failures
to distinguish observation from intervention. It cannot establish that a model
elicits the right variables from an unfamiliar bank process, measures them
correctly, or returns calibrated uncertainty after fitting to dependent data.
These are additional stages of the proposed product.

Third, a language model can propose a causal mechanism for numerical analysis.
Model Discovery Agent illustrates a separation between executable mechanism
proposals and numerical fitting, evidence approximation and experiment design.
Its finite predictive evaluations explicitly do not establish causal
identification.\cite{v12_la_2608_09696} A mechanism may predict well because
several observationally equivalent models predict equally well. The bank must
still justify which intervention the model represents and which new experiment
would distinguish economically different mechanisms.

Finally, workflow assistance concerns selection and coordination of established
operations. Causal-learn+ describes assistance with diagnosis, assumptions,
method selection and interpretation while retaining algorithm outputs and
graph ambiguity.\cite{v12_la_2606_23608} This is a useful architectural direction.
It changes the analyst's interaction with numerical tools; it does not make a
method's premises true. The following chapters keep these four contributions
separate while studying how they can work together.

\section{Text supplies a different information channel}

Write $D$ for the observed numerical records, $K$ for inspected documentary or
expert evidence, and $\mathcal M(D,K)$ for the models consistent with the
declared interpretation of both. For target $\psi$, a set-valued answer is
\[
 \Psi(D,K)=\{\psi(M):M\in\mathcal M(D,K)\}.
\]
This notation does not assert that finite data identify the population law or
that a statistical confidence procedure has been supplied. It describes which
models the analysis intends to consider. Statistical uncertainty about the law
must subsequently be propagated through the model set.

Additional reliable knowledge may exclude a model that the observed distribution
alone cannot exclude. A contract establishing that an invitation was randomized,
for example, has a different evidential role from a predictive correlation.
But a model-generated sentence repeating that randomization is common practice
does not establish what happened in the bank's historical process. The analysis
must identify the original source, the relevant population and dates, the
interpretation of its claim, and the person or measurement that supports it.

This also explains why retrieving more documents need not improve discovery.
LACR combines background knowledge, retrieved documents and optional numerical
information, and includes an UNKNOWN response. Its reported retrieval effects
differ across benchmarks, with adverse as well as favorable outcomes.
\cite{v12_la_2402_15301} Retrieval changes $K$; it does not certify $K$. Documents
may concern a different treatment definition, contain an association stated as
a cause, or duplicate the same original assertion. Ten repeated documents are
not ten independent confirmations of an edge.

An appropriate evidence object therefore retains the claim, source passage,
date, domain, units and proposed role. A passage may justify a variable
definition without justifying an arrow. A process document may establish that
one action occurs before another without excluding a common cause. A source
about another bank may suggest a hypothesis while leaving transport assumptions
unresolved. These roles determine the strength with which $K$ can restrict the
model class.

\section{Discovery methods use language judgments in different ways}

Three broad routes recur in the inspected literature: direct graph or relation
prediction, refinement of numerical discoveries, and priors or constraints used
within a search.\cite{v12_la_2402_11068} The distinction matters because it locates
where a wrong assertion enters the calculation.

In direct conditional-independence elicitation, chatPC obtains language-model
judgments from descriptions and context. Its statistical calculation concerns
the proportion of YES/NO model responses.\cite{v12_la_2406_07378} Repeated answers
can measure response variability under a prompt. Without a separate argument,
they are not repeated draws of the bank variables and do not yield a sampling
test of independence in the bank population. A highly stable wrong answer can
therefore be more dangerous than a visibly variable one.

Soft guidance provides a different failure boundary. The inspected CD-UCT
method combines language and mutual-information scores to guide graph search
without excluding otherwise valid DAGs. Language-only priors are not uniformly
beneficial in its reported comparisons.\cite{v12_la_2405_13551} Ancestral-path
constraints can likewise enter through a soft log-prior penalty, with reported
degradation under inaccurate prior information.\cite{v12_la_2306_16902} Softness
leaves room for the numerical criterion to resist the suggestion, but it does
not establish that the final graph probabilities are calibrated.

Hard restrictions can be scientifically appropriate when backed by a genuine
design fact. They are more consequential when inferred from prose: deleting a
candidate model deletes every effect value that only that model permits.
Statistical Causal Prompting converts model outputs into forced, forbidden or
unknown prior knowledge using method-specific mappings. The inspected examples
include failures as well as improvements.\cite{v12_la_2402_01454} The interface
must expose both the original assertion and the conversion rule, so a reviewer
can determine whether a numerical threshold has quietly become a scientific
exclusion.

Noisy language priors can also guide a sequential allocation of questions over
edges while retaining PAG representations.\cite{v12_la_2506_16234} This is closer
to a budgeted information problem than to asking for a finished graph in one
prompt. Its relevance to a bank still depends on whether the supplied knowledge
applies to the observed regime, whether latent structure is represented
correctly, and whether the retained ambiguity survives into the effect query.

\section{Temporal discovery is not a change of vocabulary}

A statement that collections activity affects arrears omits the action time,
available history, response horizon and possible feedback. Expanding a variable
name into a paragraph does not supply these missing indices. A temporal graph
needs explicit time labels and a rule governing contemporaneous edges, lag
limits, hidden histories, stationarity and interventions.

RealTCD combines a language-derived initial adjacency with numerical score
optimization over temporal graphs and intervention targets. Its setup includes
causal sufficiency, a known observational regime, finite lag and a DAG on
unrolled variables; its identification statement has further stated premises.
The inspected real-system assessment lacks a ground-truth DAG and evaluates
selected relationships through expert rules.\cite{v12_la_2404_14786} These
qualifications prevent a graph score from being read as a general guarantee
for a bank's latent, irregularly observed time series.

The Large Causal Models work is an adjacent numerical foundation-model route,
not simply a conversational language-model interface. The inspected model
assumes additive noise, stationarity, no latent confounders and no
contemporaneous effects, and does not detect or correct violations of these
assumptions.\cite{v12_la_2602_18662} Distinguishing the two routes avoids treating
the shared word ``model'' as a common scientific contract. One supplies semantic
proposals; the other learns a temporal mapping under a training distribution
and structural restrictions. Either can be useful within its declared regime.

For sequential servicing, a credible workflow starts with the decision history
and feasible intervention. It then asks whether the discovery output expresses
the required variables and whether a policy effect is identified over every
retained graph. An automatically generated lagged graph cannot resolve
unrecorded private information merely by assigning it a latent-node name.

\section{Interactive systems and agents add a review problem}

CausalChat combines graph suggestions, user editing and numerical updates.
Its small nonexpert study uses a fixed sequence of interfaces.
\cite{v12_la_2410_14146} Such evidence can motivate better interfaces, but it does
not isolate every learning, ordering or expertise effect. The proposed bank
evaluation needs qualified reviewers, an explicit task distribution and a
comparison that records the time and information supplied by people.

Debating agents provide another review structure. MAC uses affirmative,
negative and judge roles; its inspected comparisons also report a competitive
single-agent route for simpler graphs at lower cost.\cite{v12_la_2407_15073}
Several personas can share the same mistaken premise or model knowledge.
Agreement among them is not automatically independent corroboration. Their
usefulness should be measured through detected errors, retained correct
alternatives, decision consequences and total cost.

The oci-agent study instead emphasizes an accountable human principal over
actor/critic workflows and deterministic causal libraries. Its critic can
decline analyses with poor overlap.\cite{v12_la_2607_22443} The relevant outcome
then includes correct refusal and what remains possible for the requested
population. Reporting only completed analyses would hide precisely the
difficult cases in which an agent's judgment matters most.

An evidence-auditing agent faces a similar boundary. ARGUS audits reported
support for assumptions rather than adjudicating their truth. Its retrieval
and abstention tradeoffs and small human-label pilot leave substantial
generalization questions.\cite{v12_la_2609_30867} In banking, an omitted document,
a contradicted assumption and an assumption whose truth is unknown require
different responses. A documentation score should not become an identification
certificate.

\section{An evidence standard for the proposed banking product}

The synthesis supports a modular research program. Language assistance proposes
questions, variables, graph alternatives and evidence links. Formal operations
check declared structures and calculate identified functionals or sets in
supported regimes. Numerical procedures estimate those objects with their own
sampling and approximation arguments. Human and institutional evidence supply
facts about assignment, ownership and feasibility that neither interface can
manufacture.

Each interface needs an explicit refusal or unknown state. A retrieved passage
with the wrong unit should fail target review. A graph class outside a backend's
domain should return unsupported. A valid but uninformative effect set should
remain visible. A conflict between textual and numerical evidence should retain
both sources for investigation instead of selecting whichever yields a cleaner
answer.

CausalDS makes narrated mechanisms, observations and private scoring objects
separate parts of a benchmark, and its results distinguish symbolic answers,
uncertainty, abstention and cost.\cite{v12_la_2607_08093} The proposed evaluation
adopts that separation as a design principle. It does not borrow a published
model score as evidence for the local prototype. The next chapters develop the
translation, diagram, discovery, workflow and economic consequences of this
architecture; its empirical value still requires independent mechanisms and
an actual declared backend.
